% ============================================================================
%  Chapitre 1 — Introduction
%  Cible : 4 à 5 pages. À RÉDIGER EN DERNIER (J12), après le chapitre 6.
% ============================================================================
\chapter{Introduction}
\label{chap:introduction}

\section{Contexte général}
\label{sec:intro-contexte}

\aecrire{4 à 6 paragraphes. Entonnoir : (1) électrification et numérisation des
réseaux électriques, place du comptage d'énergie communicant ; (2) le produit
n'est plus un équipement figé mais une plateforme logicielle mise à jour sur
le terrain ; (3) conséquence : la qualité et la sécurité du logiciel embarqué
deviennent un enjeu produit, pas un enjeu d'atelier ; (4) l'industrialisation
de la chaîne de production logicielle (CI/CD) devient donc un sujet
d'ingénierie à part entière.}

\section{Problématique}
\label{sec:intro-problematique}

Le développement du firmware d'un produit industriel embarqué cumule des
contraintes que l'on rencontre rarement simultanément dans le développement
logiciel classique :

\begin{itemize}[nosep]
  \item des temps de construction très longs, la génération d'une distribution
        Linux embarquée complète se comptant en heures et non en minutes ;
  \item un environnement de compilation complexe et fragile, associant chaînes
        de compilation croisée, bibliothèques système et outils propriétaires ;
  \item des exigences réglementaires et normatives de cybersécurité
        (IEC~62443-4-1, règlement (UE) 2024/2847 dit \emph{Cyber Resilience
        Act}) imposant la traçabilité des composants et la détection précoce
        des vulnérabilités ;
  \item une exécution des traitements sur une infrastructure mutualisée et
        éphémère, où aucun état ne persiste d'une exécution à l'autre.
\end{itemize}

\medskip
\noindent
Ces contraintes conduisent à la question de recherche suivante :

\begin{quote}
\itshape
Comment concevoir et mettre en œuvre une chaîne d'intégration et de livraison
continues reproductible, sécurisée et mesurable pour le firmware d'un produit
industriel embarqué, sous contraintes de temps de construction longs,
d'exigences normatives de cybersécurité et d'infrastructure d'exécution
éphémère ?
\end{quote}

\noindent
Cette question se décline en quatre sous-questions traitées respectivement dans
les sections~\ref{sec:impl-conteneur} à~\ref{sec:impl-observabilite} :

\begin{description}[style=nextline,leftmargin=1.5cm]
  \item[QR1 — Reproductibilité] Comment garantir qu'une même révision du code
        source produise un artefact identique, indépendamment de la machine et
        de la date d'exécution ?
  \item[QR2 — Sécurité] Comment intégrer l'analyse statique du code et
        l'analyse de composition logicielle dans le cycle de développement sans
        dégrader la vélocité des équipes de développement ?
  \item[QR3 — Mesurabilité] Quels indicateurs instrumenter pour piloter la
        qualité de la chaîne, et comment les collecter automatiquement ?
  \item[QR4 — Passage à l'échelle] Comment dimensionner et administrer
        l'infrastructure d'exécution de manière automatisée et maîtrisée en
        coût ?
\end{description}

\section{Objectifs et périmètre}
\label{sec:intro-perimetre}

\aecrire{Formuler 4 à 5 objectifs opérationnels mesurables, puis délimiter
explicitement ce qui est hors périmètre.}

Le périmètre de ce travail est la \emph{chaîne de production} du firmware, et
non le firmware lui-même. Sont explicitement exclus : la conception
fonctionnelle de l'application embarquée, le déploiement des mises à jour sur
le parc installé, ainsi que la conception matérielle du produit au-delà des
contraintes qu'elle impose à la compilation.

\section{Contributions}
\label{sec:intro-contributions}

Les contributions de ce travail sont les suivantes :

\begin{enumerate}[nosep]
  \item la conteneurisation de l'environnement de compilation applicatif, avec
        épinglage des versions de l'ensemble de la chaîne d'outils
        (chapitre~\ref{chap:mise-en-oeuvre}, section~\ref{sec:impl-conteneur}) ;
  \item la conception de deux chaînes d'intégration distinctes et
        complémentaires, l'une pour la couche applicative et l'autre pour la
        distribution Linux embarquée (sections~\ref{sec:impl-pipeline-app}
        et~\ref{sec:impl-pipeline-yocto}) ;
  \item l'intégration automatisée de l'analyse statique et de l'analyse de
        composition logicielle avec restitution des résultats au développeur
        (section~\ref{sec:impl-securite}) ;
  \item l'automatisation de l'infrastructure d'exécution par
        \emph{infrastructure as code} (section~\ref{sec:impl-runners}) ;
  \item la mise en place d'une chaîne d'observabilité restituant l'évolution
        des indicateurs de qualité dans le temps
        (section~\ref{sec:impl-observabilite}) ;
  \item une évaluation quantitative de la chaîne obtenue
        (chapitre~\ref{chap:evaluation}).
\end{enumerate}

\section{Méthodologie}
\label{sec:intro-methodologie}

\aecrire{Assumer explicitement la posture : étude de cas industrielle unique,
de type recherche-action. Décrire le déroulement : analyse de l'existant,
recueil des exigences auprès des équipes, conception itérative, mise en
production incrémentale, mesure. Citer la littérature sur l'étude de cas en
génie logiciel (Runeson \& Höst) pour asseoir la démarche.}

\section{Organisation du mémoire}
\label{sec:intro-plan}

Le chapitre~\ref{chap:contexte} présente le contexte industriel, le produit et
l'état initial de la chaîne de développement. Le chapitre~\ref{chap:etat-art}
dresse l'état de l'art de l'intégration continue, de son application aux
systèmes embarqués et des pratiques de sécurisation de la chaîne
d'approvisionnement logicielle. Le chapitre~\ref{chap:besoins} formalise les
exigences et présente l'architecture retenue ainsi que la justification des
choix technologiques. Le chapitre~\ref{chap:mise-en-oeuvre} détaille la mise en
œuvre. Le chapitre~\ref{chap:evaluation} évalue quantitativement les résultats
obtenus et en discute les limites. Le chapitre~\ref{chap:conclusion} conclut et
ouvre sur les perspectives.
